\section{Sound characterisation of the Smyth probabilistic must preorder}\label{app:smyth}\label{sec:must}

\paragraph{Observers and tested systems.} An \emph{observer} is a program $T$
in which the guard $\mathbf{1}$ marks success: $T$ is \emph{successful} when
$\mathbf{1}$ occurs at top level, through parallel compositions,
restrictions, conditionals whose guard evaluates and sums (so that success
is invariant under $\equiv$). A program $P$ with a memory $\rho$ is tested
by running $\conf{\rho}{P \parallel T}$; the observer has no memory of its
own, its prefixes $U(\tilde q)$, $\M(\tilde q \rhd x)$ and
$\mathsf{init}(\tilde q)$ act on the shared memory. Following
\cite{Ceragioli27} we resolve non-determinism by schedulers, with a lifting
to distributions that keeps every branch of a measurement under the same
choices. A \emph{distribution} is a formal sum $\sum_i \rho_i \cdot R_i$ of
programs weighted by partial density operators; its \emph{mass} is
$\sum_i \tr(\rho_i)$, its \emph{successful mass} the mass of its successful
components. The \emph{scheduled step} of a configuration,
$\rho \cdot R \lstep{s} \Delta$, is a silent step of $\conf{\rho}{R}$ with
choice $s$ (\cref{fig:conf}) in which the outcomes of a measurement are
grouped: $\Delta = \sum_m \br \cdot R_m$ over the outcomes $m$ of non-zero
probability for \textsc{Meas} (a distribution is its support, as in a
pLTS), and $\rho' \cdot R'$ for the other rules. A distribution moves \emph{as a whole}
under one choice: $\Delta \lstep{s} \Delta'$ when at least one component
moves, every non-successful component that can do $s$ does it, the
successful ones stay, and the non-successful ones that cannot do $s$ are
\emph{lifted}: they stay (lifting Stay) or they disappear (lifting Drop).
A \emph{run} is an infinite sequence $\Delta_0 \lstep{s_1} \Delta_1
\lstep{s_2} \cdots$ that stutters once no choice moves any component; its
\emph{value} is the supremum of the successful masses of the $\Delta_i$.

\begin{definition}[Smyth must]\label{def:smyth}
For a class $\mathcal{O}$ of observers, $(P, \rho) {\smyth}_{\mathcal{O}} (Q, \sigma)$
iff for every $T \in \mathcal{O}$ and every run of $\rho \cdot (Q \parallel T)$
there is a run of $\sigma \cdot (P \parallel T)$ of value at most the value
of the former.
\qed
\end{definition}

This is the must preorder of Deng et al.~\cite{DvGHM08} with the
resolutions restricted to one word of choices for the whole distribution,
as in the definition of Ceragioli et al.~\cite[Definition 7]{Ceragioli27}, whose resolutions are the
probabilistic runs of a scheduler. The definition is generic (the Smyth
must of \texttt{pMust\_smyth.v}, for any scheduled probabilistic system
and either lifting); lqCCS only provides the step and the composition. We
write $\smyth$ for the lifting Stay and $\smyth^{\mathsf{D}}$ for Drop.
With Stay, $\ite{x = 0}{t_1{:}\tau.P}{t_2{:}\tau.Q}$ after a measurement
completes both branches under the word $t_1 t_2$; with Drop, the word $t_1$
discards the branch $x = 1$, which counts as failed, as in
\cite{Ceragioli27} where a path that cannot follow the scheduler has no
run. Drop extends that definition in two respects: successful components
are frozen instead of disappearing when another component moves on, and
the runs are infinite, so that a divergent process is the bottom of both
preorders ($\Omega = \mathsf{rec}\,X.\,t{:}\tau.X$ satisfies $(\Omega,
\rho) \smyth (Q, \rho)$ and $(\Omega, \rho) \smyth^{\mathsf{D}} (Q, \rho)$
for every $Q$, \texttt{Omega\_bottom}), whereas finite maximal schedulers make
it incomparable. The two preorders differ: for
$q = \M(q \rhd x).(\ite{x=0}{t_1{:}\tau.a!0}{t_2{:}\tau.a!0})$ and
$q' = \M(q \rhd x).(\ite{x=0}{t{:}\tau.a!0}{t{:}\tau.a!0})$, the observer
$a?z.\mathbf{1}$ gives the value $1$ to both under Stay, while under Drop
$q$ has the value $1/2$ and $q'$ the value $1$, so that
$q \smyth^{\mathsf{D}} q'$ strictly.

\paragraph{The class of observers.} We characterise ${\smyth}_{\mathcal{O}}$
for the observers without internal choice: no $\tau$ prefix, no recursion,
no restriction, no channel used both for input and for output (no
communication inside the observer), $\mathbf{1}$ never in a sum with a
prefix (success is persistent), each tag naming one prefix, and tags
disjoint from those of the tested programs, as in the work of Ceragioli et al.~\cite{Ceragioli27}.
For such an observer the offers at a point of the interaction only depend on
the \emph{co-trace} that led there, the sequence of its co-actions (tag of
the observer, action, with the value received or the outcome measured).

\begin{remark}
An observer with a $\tau$ choice can exploit \emph{timing}. Let
$P = \M(q \rhd x).c!x.(\ite{x=0}{a!0}{b!0})$ and
$Q = \M(q \rhd x).(\ite{x=0}{c!x.a!0}{t_5{:}\tau.c!x.b!0})$, and
$T = c?y.(t{:}\tau.a?z.\mathbf{1} + t'{:}\tau.b?z.\mathbf{1})$. For $Q$
the scheduler delivers the branch $x=0$ first, chooses $t'$ for it, then
$t_5$, delivers the branch $x=1$ and chooses $t$: no branch succeeds. For $P$
both branches reach $T$ at the same step and $t$ or $t'$ applies to both:
value $1/2$ always. So $P \not\smyth Q$ although $P$ and $Q$ only differ by
a $\tau$ in one branch, and the preorder is not insensitive to $\tau$. The
scheduler dates the components through the tags that only one branch has,
and the choice of the observer may depend on that date. Characterising this
preorder would need co-traces with dates; we leave it out.
\end{remark}

\paragraph{The LTS $\DD$.} The alternative preorder is defined on the
LTS $\DD$ of distributions of \emph{components} $(n, (P, C, E))$ weighted
by a memory, where $n$ is the co-trace of the component, $E$ the outputs
(tag, channel, value) the observer offers at $n$, recorded when $n$ is
reached, and $C$ those of them the program has consumed. Outputs are
asynchronous: the observer does not see when they are consumed, so a
consumption is a silent step of the component, which does not extend the
co-trace; $C$ only forbids a second consumption. The labels of $\DD$ are
the choices $s$ of the program (a consumption being the choice $(t,k)$ of
the tags involved) and the \emph{instructions} $(k, t)$, ``the tag $k$ of
the observer, answered by the tag $t$ of the program''. Three parameters
describe the observer, its \emph{shape}: its \emph{menu} $I(n)$, the inputs
$k{:}c?$, operations $k{:}U(\tilde q)$, measurements $k{:}\M(\tilde q)$ and
initialisations it offers at the co-trace $n$; its \emph{outputs} $O(n)$ at
$n$; and its \emph{success co-traces} $\succn(n)$, closed under extension.
\cref{tab:dlts} gives the step of one component; a distribution moves as
above, the components whose co-trace is in $\succn$ being frozen.

\begin{table}[ht]
\centering
\footnotesize
$\begin{array}{l@{\quad}rcl@{\quad}l}
\textsc{Ch} & \rho \cdot (n,(P,C,E)) &\lstep{s}& \textstyle\sum_i \rho_i \cdot (n,(P_i,C,E))
  & \rho \cdot P \lstep{s} \sum_i \rho_i \cdot P_i \\[2pt]
\textsc{Recv} & \rho \cdot (n,(P,C,E)) &\lstep{(t,k)}& \rho \cdot (n,(P',C \cdot (k,c,v),E))
  & (k,c,v) \in E,\ k \notin C,\ P \step{c?v}{t} P' \\[2pt]
\textsc{Send} & \rho \cdot (n,(P,C,E)) &\lstep{(k,t)}& \rho \cdot (n',(P',C,O(n')))
  & k{:}c? \in I(n),\ P \step{c!v}{t} P',\ n' = n \cdot k{:}c?v \\[2pt]
\textsc{Op} & \rho \cdot (n,(P,C,E)) &\lstep{(k,t)}& \upd{U} \cdot (n',(P,C,O(n')))
  & k{:}U(\tilde q) \in I(n),\ n' = n \cdot k{:}U(\tilde q) \\[2pt]
\textsc{Init} & \rho \cdot (n,(P,C,E)) &\lstep{(k,t)}& \inirho \cdot (n',(P,C,O(n')))
  & k{:}\mathsf{init}(\tilde q) \in I(n),\ n' = n \cdot k{:}\mathsf{init}(\tilde q) \\[2pt]
\textsc{Meas} & \rho \cdot (n,(P,C,E)) &\lstep{(k,t)}& \textstyle\sum_{m,\ \br \neq 0} \br \cdot (n_m,(P,C,O(n_m)))
  & k{:}\M(\tilde q) \in I(n),\ n_m = n \cdot k{:}\M(\tilde q)!m
\end{array}$
\caption{The step of a component of $\DD$, for a shape $(I, O, \succn)$.}
\label{tab:dlts}
\end{table}

\begin{definition}[Alternative preorder]\label{def:alt}
$(P,\rho) \alt (Q,\sigma)$ iff for every shape $(I, O, \succn)$ with
$\succn$ closed under extension, for every run of
$\sigma \cdot ([],(Q,[],O([])))$ in $\DD$ there is a run of
$\rho \cdot ([],(P,[],O([])))$ in $\DD$ of value at most the value of the
former, the value being the mass of the components whose co-trace is in
$\succn$.
\qed
\end{definition}

The runs of $\DD$ are the trees of runs of the configuration LTS following
one word of choices and instructions: the co-traces of the components are
the co-traces of $\conf{\rho}{P}$ in the LTS $\rho \lpar P$ with the actions
of the memory exposed to the observer (rule Left of \cref{def:par}), and
$\DD$ keeps them together because a word applies to all of them. A run
stops only when every unfrozen component is a \emph{waiting state}, with no
silent step and no offer of $I(n)$ it can answer: this is where the
acceptance sets of the classical characterisation live, $\coR(P') \cap I(n)
= \emptyset$ for every leaf $(n, P')$. Taking one co-trace at a time is not
possible: two processes that, after a measurement, send the outcome and
choose $a$ or $b$ with the same tags in both branches (the choice correlated
or anti-correlated with the outcome) have the same leaves co-trace by
co-trace and are told apart by $c?y.(\ite{y=0}{a?z.\mathbf{1}}{b?z.\mathbf{1}})$.
The quantum state only enters through the masses $\tr(\rho')$ of the leaves:
the operations and measurements of the observer are instructions inside the
run, so no Löwner comparison of the leaves is needed, and it would be
incomplete (a program that sends $|0\rangle$ or $|1\rangle$ by an internal
choice is below one that sends $I/2$).

\paragraph{The theorems.} The shape of an observer $T$ is the triple
$(\menu_T, O_T, \succn_T)$: $\menu_T(n)$ and $O_T(n)$ are the offers and
the outputs of the state of $T$ after the co-trace $n$, and $\succn_T(n)$
holds when that state, or the state at a prefix of $n$, is successful. Write $\mathcal{O}_{P,Q}$ for the
observers of the class above whose tags are apart from those of $P$ and
$Q$, and $\alt_{\mathcal{O}}$ for \cref{def:alt} restricted to the shapes of
the observers of $\mathcal{O}_{P,Q}$.

\begin{theorem}[Soundness and characterisation]\label{thm:char}
Let $P, Q$ be programs without $\mathbf{1}$, and let $\DD$, $\alt$ and
$\smyth$ be taken with the same lifting, Stay or Drop. Then
\[
(P,\rho) \alt (Q,\sigma)
\ \Longrightarrow\
(P,\rho) \alt_{\mathcal{O}} (Q,\sigma)
\ \Longleftrightarrow\
(P,\rho) {\smyth}_{\mathcal{O}_{P,Q}} (Q,\sigma).
\]
\end{theorem}

\begin{proof}
Fix an observer $T$. A component $\rho \cdot (P \parallel T')$ of the tested
system is related to a component $\rho \cdot (n,(P,C,O_T(n)))$ of $\DD$
when $T'$ is the state of $T$ after the co-trace $n$ followed by the
outputs $C$. Related distributions have the same successful mass ($P$ has
no $\mathbf{1}$, and the outputs of $C$ do not change the success of $T'$).
Every scheduled step of one is a scheduled step of the other with the
same successors, related: a choice of the program is the same choice on
both sides, the consumption of an output of $T'$ included (\textsc{Recv});
an input of $T'$ served by $P$ is the instruction $(k,t)$ of the tags
involved (\textsc{Send}); an operation of $T'$ on the memory is the
instruction $(k,t)$ of its tag (\textsc{Op}, \textsc{Init}, \textsc{Meas});
the tags being disjoint and unique, the choice determines the instruction
and conversely. The
consumptions commute with the other co-actions of the observer and keep its
other offers, so the menu at the co-trace $n$ is what $T'$ offers. A
distribution is then stable iff the related one is, the lifting of a
component is the lifting of the related one (both stay, or both are
successful or both disappear), and a run of one system is mirrored by a run
of the other with the same successful masses, which needs dependent choice
to build the mirrored run. The two implications of
the equivalence apply this transfer to the run of $Q$, then to the witness
run of $P$; the first implication is the instance of $\alt$ at the shape of
$T$.
\qed
\end{proof}

\paragraph{The almost-sure must.} The same transfer characterises the
must with probability one. $P$ passes $T$ almost surely when every run of
$\rho \cdot (P \parallel T)$ has the value $\tr(\rho)$, its successful mass
tending to its mass, and $(P,\rho) \sqsubseteq_{\mathsf{as}, \mathcal{O}}
(Q,\sigma)$ when $Q$ passes almost surely every $T \in \mathcal{O}$ that
$P$ passes almost surely. Let $(P,\rho) \alt_{\mathsf{as}} (Q,\sigma)$ iff
for every shape $(I, O, \succn)$ with $\succn$ closed under extension, if
every run of $\rho \cdot ([],(P,[],O([])))$ in $\DD$ has the value
$\tr(\rho)$ then every run of $\sigma \cdot ([],(Q,[],O([])))$ has the value
$\tr(\sigma)$. Under the hypotheses of \cref{thm:char}, and with the same
lifting on both sides,
\[
(P,\rho) \alt_{\mathsf{as}} (Q,\sigma)
\ \Longrightarrow\
(P,\rho) \alt_{\mathsf{as}, \mathcal{O}} (Q,\sigma)
\ \Longleftrightarrow\
(P,\rho) \sqsubseteq_{\mathsf{as}, \mathcal{O}_{P,Q}} (Q,\sigma),
\]
since a run and its mirrored run have the same successful masses and the
same mass (\texttt{alt\_as\_sound}, \texttt{alt\_as\_obs\_iff}). Under Drop
a run has full value only if no component ever disappears. The Smyth must
is finer than the almost-sure must: if below every run of $Q \parallel T$
there is a run of $P \parallel T$ and the latter all have full value, so
do the former (\texttt{ctx\_pre\_sm\_as}).

\begin{remark}
\cref{def:alt} quantifies over all menus, including menus that are the shape
of no observer (infinitely many offers, a tag offering different kinds of
actions at different co-traces, menus not uniform in the value received).
The equivalence is therefore stated for the shapes of observers. The
quantification over the menu is not removed as in the classical
characterisation, where the witness $p'$ with $\coR(p') \subseteq
\coR(q')$ refuses every menu that $q'$ refuses: here the witness is a run,
and the menu also drives it afterwards, an offer at a waiting state forcing
it to move. The \emph{uniform} form, where the witness run of $P$ is chosen
before the menu and is a run under every menu under which the run of $Q$ is
one (\texttt{alt\_pre\_unif}: the witness uses only the instructions of the
run of $Q$, stops only where that run stops, and its leaves refuse what the
leaves of $Q$ at the same co-trace refuse), implies $\alt$; so does the
weaker \emph{scheduled-word} form, where only the sequence of items and the
choices of the components are fixed before the menu, an item that no
component can take being skipped, and the run has to be maximal under every
menu (\texttt{alt\_pre\_det}, \texttt{alt\_pre\_det\_alt\_pre}). Whether
either converse holds is not settled, for either lifting. The
difficulty is not the one of the classical theory: the witness may have to
park a component at a co-trace that the run of $Q$ never visits, and what
it does there is decided by the menu; but the shapes punish such a
component, since an offer at that co-trace can lead to success for $P$
alone, and every candidate counterexample we examined (found by an
exhaustive search over small programs and bounded menus) was eventually
separated by a deeper shape. It would fail, for both liftings, if the
outcomes of probability zero of a measurement were kept in the
distributions: a component $d!v$ of
mass zero can neither stop under the maximal menu nor move under the
minimal one, while it costs nothing to $\alt$. This is why the scheduled
step only keeps the outcomes of non-zero probability; from a non-zero
memory, every component of a run then has a non-zero memory
(\texttt{drun\_nonzero}).
\end{remark}
